\section{The Architect's Question}\label{the-architects-question}

As the ``From Token to Fleet'' handbook reaches its final chapters, we
turn from technical mechanisms to the governance of those mechanisms.
With \textasciitilde2,000 users, a RAG Q\&A system, and a 70B dense FP16
deployment, the question becomes: how does the team capture, justify,
and evolve the architectural choices that keep the fleet running? The
Architecture Decision Record (ADR) is the answer --- a lightweight,
text-first practice that records not just what was decided, but why,
what alternatives were considered, and what trade-offs were accepted.
This chapter explores the ADR pattern in the context of a production AI
system, providing a reusable template and a worked example centered on
model-shape choices.

\section{Concept}\label{concept}

An Architecture Decision Record is a single Markdown file that captures
a significant architectural decision and its context. The format was
popularized by Michael Nygard and has been adopted by teams ranging from
startups to enterprises managing complex AI/ML stacks. Each ADR follows
a minimal structure:

\begin{itemize}
\tightlist
\item
  \textbf{Title} and status (proposed, accepted, superseded, deprecated)
\item
  \textbf{Context} --- the problem, constraints, and goals that prompted
  the decision
\item
  \textbf{Decision} --- the chosen approach, concise and explicit
\item
  \textbf{Consequences} --- outcomes, both positive and negative, that
  flow from the decision
\item
  \textbf{Alternatives considered} --- other paths rejected, with brief
  rationale for rejection
\end{itemize}

The key distinction between an ADR and a standard design document is
intent: an ADR is meant to be living, reviewed in pull requests, and
occasionally superseded. It answers the ``why'' at the moment the
decision was made, preserving it for future engineers who must reason
about the system years later.

In a RAG/Q\&A fleet with 70B models, ADRs prevent the ``why did we
choose this tokenizer?'' or ``why FP16 over BF16?'' questions from
becoming tribal knowledge. They also integrate naturally with git ---
each ADR lives in version control, linked to the commit or PR that
implemented the decision.

\section{Mental Model}\label{mental-model}

The mental model underpinning ADRs is simple: architectural decisions
are irreversible (or costly to reverse) choices that deserve permanent
documentation. Unlike code comments, which explain what a function does,
ADRs explain why the system is structured as it is. This mental model
has three layers:

\begin{enumerate}
\def\labelenumi{\arabic{enumi}.}
\tightlist
\item
  \textbf{Problem awareness} --- What forces are at play? (load
  patterns, latency requirements, budget constraints, team expertise)
\item
  \textbf{Choice architecture} --- What design options exist? What are
  their properties?
\item
  \textbf{Outcome acceptance} --- What will we live with? What will we
  monitor?
\end{enumerate}

For AI/ML systems, this model extends to decisions about model
quantization, retrieval chunk sizes, embedding dimensions, and
infrastructure trade-offs. Each decision carries FACT/DERIVED/HYPOTHESIS
evidence tags, making the ADR itself a source of verifiable claims
rather than opinion.

The ADR format also enforces a habit: before committing to a decision,
the team must explicitly articulate alternatives. This alone often
reveals overlooked options or clarifies why the chosen path is genuinely
optimal.

\section{Worked Example: Model-Shape
ADR}\label{worked-example-model-shape-adr}

\textbf{Title:} ADR 0016 --- Model Quantization: FP16 vs BF16 for 70B
Inference

\textbf{Status:} Accepted

\textbf{Context:} The fleet runs a 70-billion-parameter dense model for
RAG Q\&A serving. Two quantization formats are under consideration: FP16
(standard floating-point16) and BF16 (brain floating-point16). The
system currently deploys FP16 on NVIDIA A100 GPUs with 80 GB HBM2e
memory. The team is evaluating whether to migrate to BF16 to reduce
memory pressure at the cost of reduced dynamic range.

\textbf{Decision:} Continue with FP16 quantization for the 70B model.
BF16 migration is deferred to a future hardware refresh.

\textbf{Rationale:} - FP16 provides sufficient precision for Q\&A use
cases; empirical evaluation on the retrieval and generation pipelines
shows no measurable drop in exact-match scores. - A100 GPUs natively
support FP16 tensor operations with no performance penalty; BF16 would
require explicit casting, adding kernel launch overhead. - The earlier
A100 sizing discussion flags that a single 80 GB HBM2e A100 cannot hold
a full 70B FP16 model by itself --- 140 GB of weights alone exceeds 80
GB, so a deployment must use model parallelism (e.g., 2×80 GB A100 for
the \textasciitilde140 GB weights with KV in a portion of the pool) or
quantization (FP8/8-bit) to fit a single card. The ADR therefore records
the real constraint: FP16 residency is \textasciitilde164 GB with 9.2K
KV on the canonical 70B (\hyperref[chap:7]{Chapter 7}), requiring multiple GPUs or
compression. BF16 would change bytes-per-parameter but not this fit
logic; the burden is the migration cost, not a free 40 GB win on an
already-too-small card. - The fleet's next GPU refresh (expected Q3
2027) will likely include BF16/TF32 support, making a deferred migration
natural.

\textbf{Consequences:} - \textbf{Positive:} No code change required;
zero risk of introducing inference bugs; existing monitoring and
alerting remain valid. - \textbf{Negative:} \textasciitilde40 GB more
memory per model instance compared to BF16; if GPU prices or
availability shift, the fleet may need to carry fewer concurrent models.
- \textbf{Monitoring:} Track per-request latency, GPU memory
utilization, and model output quality on a held-out Q\&A set. If quality
degrades, reassess quantization choices.

\textbf{Alternatives Considered:} 1. \textbf{Pure FP32:} Rejected ---
280 GB per model instance, exceeding A100 memory by 200+ GB. Not viable
without model parallelism. 2. \textbf{8-bit Int8 quantization:} Rejected
--- while memory reduces to \textasciitilde70 GB, accuracy drops on
multi-step reasoning tasks; not acceptable for the Q\&A SLA. 3.
\textbf{Model parallelism (pipeline parallel):} Rejected --- introduces
cross-GPU communication latency; degrades user-perceived response time
for short queries.

\textbf{Evidence Tags:} FACT: A100 FP16 tensor core throughput = 312
TFLOPS. DERIVED: 70B FP16 model footprint ≈ 140 GB parameters + 30 GB
activations ≈ 170 GB total. HYPOTHESIS: BF16 migration would reduce
memory by \textasciitilde40\% with \textless0.5\% quality impact, based
on {[}Schema et al., 2023{]}.

\subsection{3.1 A copy-ready blank
template}\label{a-copy-ready-blank-template}

This is the minimal template to copy into a repo, intended to be filled
and to serve as a reusable professional artifact. Every section maps to
a discipline from this handbook; the model-shape ADR above is one filled
example, and \hyperref[chap:22]{Chapter 22}'s loop supplies the architecture-decision
context.

\begin{verbatim}
# ADR-NNNN — <Short Decision Title>

## Status
<proposed | accepted | superseded | deprecated>   (date)

## Decision (the committed choice, one paragraph, reversible in principle)

## Architecture Decision Context   (from \hyperref[chap:22]{Chapter 22}'s loop)
- Business objective:  <what outcome this enables>
- Workload:            <characterization: tokens, rps, peak>
- Requirements+SLOs:   <quality / latency / availability, measurable>
- Constraints:         <privacy, region, ops, procurement, budget>

## Decision

## Alternatives considered
1. <option> — <why rejected>
2. <option> — <why rejected>

## Evidence
- <claim> — FACT/DERIVED/HYPOTHESIS [1P]/[2°]/[VERIFY]
- Unit check:  <does each derived quantity's unit make sense?>
- Sanity check: <is each result physically possible? (e.g. MFU ≤ 1)>

## Trade-offs & consequences
- Pro:   <expected positive effects>
- Con:   <expected negative effects>
- Risk:  <residual uncertainty>

## Decision owner
<person / team>

## Validation plan
<how this decision will be benchmarked/verified post-deploy>

## Rollback plan
<how we unwind if the decision fails its SLO>

## Review trigger
<what event (metric, date, new evidence) re-opens this ADR>
\end{verbatim}

The three blocks most often skipped are the \textbf{Architecture
Decision Context}, the \textbf{unit check}, and the \textbf{decision
owner + review trigger} --- and skipping exactly those three is what
turns an ADR into a decision log instead of a decision \emph{apparatus}.
The context ties the technical choice to a business outcome; the unit
check enforces \hyperref[chap:22]{Chapter 22}'s evidence discipline; the owner and review
trigger make it a living document rather than a tombstone.

\section{Measurement}\label{measurement}

ADRs are only as good as the evidence they embed. This chapter
recommends that each ADR include at least two derived quantities ---
quantities computed from raw data, not merely observed. In the
model-shape ADR above, the derived quantities are:

\begin{enumerate}
\def\labelenumi{\arabic{enumi}.}
\tightlist
\item
  \textbf{70B FP16 model footprint} ---
  \(W = N \times \text{bytes-per-param} = 70 \times 10^9 \times 2 \text{ B} = 140\)
  GB parameters, plus \textasciitilde30 GB activations (typical
  profiling-derived overhead) ≈ \(170\) GB total. Derived from the model
  parameter count and a typical activation overhead factor estimated
  from profiling runs.
\item
  \textbf{BF16 migration would reduce memory by \textasciitilde40\% with
  \textless0.5\% quality impact} --- the parameter footprint scales
  \textasciitilde linearly with byte-width: FP16 (\(2 \text{ B/param}\))
  vs BF16-to-FP32 mixes, giving roughly a \(2\times\) byte-width ratio
  that translates to a \textasciitilde40\% memory reduction for the same
  model; validated by a small-scale accuracy benchmark on 1\% of the
  validation set.
\end{enumerate}

Additional measurable quantities that should be tracked (even if not all
included in the ADR itself) include:

\begin{itemize}
\tightlist
\item
  GPU memory utilization per request (derived from NVIDIA management
  library stats)
\item
  Per-request latency p99 (derived from request-timing histograms)
\item
  Exact-match score on the Q\&A dev set (derived from evaluation harness
  runs)
\end{itemize}

Each derived quantity should be tagged with its evidence source
{[}1P{]}/{[}2°{]}/{[}VERIFY{]}, linking to a profiling script, a
benchmark run, or a verified observation. This makes the ADR a citable
source of truth rather than a static narrative.

\section{Common Mistakes}\label{common-mistakes}

When adopting ADRs, teams frequently fall into patterns that undermine
the format's purpose. Here are the most common, with fixes:

\textbf{\hyperref[tab:25.1]{Table 25-1}} --- Common ADR adoption mistakes and their fixes.

\begin{longtable}[]{@{}
  >{\raggedright\arraybackslash}p{0.5000\linewidth}
  >{\raggedright\arraybackslash}p{0.5000\linewidth}@{}}
\toprule\noalign{}
\begin{minipage}[b]{\linewidth}\raggedright
Mistake
\end{minipage} & \begin{minipage}[b]{\linewidth}\raggedright
Fix
\end{minipage} \\
\midrule\noalign{}
\endhead
\bottomrule\noalign{}
\endlastfoot
\textbf{Writing ADRs after the fact} --- retroactive documentation loses
the context of the decision trade-offs. & ADRs are written as part of
the decision process, ideally in the same PR that implements the
change. \\
\textbf{Too much detail} --- documenting every incremental choice
clutters the record and discourages reading. & Limit ADRs to decisions
that affect system behavior, architecture, or cost. Routine choices
(library upgrades without API change) need no ADR. \\
\textbf{No status tracking} --- an ADR that is never updated becomes
stale and erodes trust. & Use status fields (proposed → accepted →
superseded) and treat the ADR as a living document reviewed in PRs. \\
\textbf{Missing alternatives} --- without explicit alternatives, the ADR
reads like a proclamation, not a decision log. & Always include an
``Alternatives considered'' section, even if the list is short. \\
\textbf{No evidence links} --- claims without references become
folklore. & Every derived quantity and key claim should reference a data
source, script, or benchmark run. \\
\textbf{Treating ADRs as permanent} --- some decisions do change. & If a
decision is superseded, update the ADR status and add a superseding ADR
link. Do not delete history. \\

\label{tab:25.1}\end{longtable}

\section{Architecture Consequence}\label{architecture-consequence}

Introducing ADRs into the ``From Token to Fleet'' handbook has several
architecture-level consequences:

\begin{itemize}
\tightlist
\item
  \textbf{Decision provenance:} Every significant choice in the fleet's
  evolution is traceable to a git commit and a Markdown file. New
  engineers can audit why the system is as it is, without relying on
  Slack history or outdated wikis.
\item
  \textbf{RQO (Reasoned Quality Optimization):} By forcing the
  articulation of alternatives and evidence, ADRs make the trade-off
  surface explicit. Teams can point to an ADR and say ``we chose FP16
  over BF16 because of X, Y, Z'' --- and those reasons are verifiable.
\item
  \textbf{Reduced cognitive load:} Engineers no longer need to re-derive
  the rationale for every architectural choice. The ADR serves as a
  single source of truth.
\item
  \textbf{ADR proliferation:} As the fleet grows, the number of ADRs
  will grow. A naming convention (ADR 0001, ADR 0002, \ldots) and a
  top-level index (ADR index.md) prevent the record from becoming
  unwieldy.
\item
  \textbf{Integration with CI:} ADRs can be validated as part of the CI
  pipeline --- e.g., ensuring every ADR has a status, a rationale, and
  at least one evidence tag. This automation reinforces the habit.
\end{itemize}

On the negative side, ADRs add a small writing overhead. For a team of
\textasciitilde10 engineers making \textasciitilde1--2 significant
architectural decisions per month, this is a manageable cost. The payoff
--- reduced on-call noise, faster onboarding, and fewer ``why did we do
it this way?'' debates --- more than compensates.

\begin{figure}
\centering
\pandocbounded{\includegraphics[keepaspectratio,alt={\hyperref[fig:25.1]{Fig 25.1} --- The Anatomy of an Architecture Decision Record {[}ILLUSTRATIVE conceptual{]}},width=\maxwidth{\textwidth},keepaspectratio]{figures/fig-25-2501.pdf}}
\caption{\hyperref[fig:25.1]{Fig 25.1} --- The Anatomy of an Architecture Decision Record
{[}ILLUSTRATIVE conceptual{]}}

\label{fig:25.1}\end{figure}

\section{What We Still Don't Know}\label{what-we-still-dont-know}

Despite the structured format, several questions remain open and would
benefit from future ADRs or empirical studies:

\begin{enumerate}
\def\labelenumi{\arabic{enumi}.}
\tightlist
\item
  \textbf{ADR fatigue:} Does the presence of many ADRs overwhelm the
  team, leading to superficial writing or skipped documentation?
\item
  \textbf{Evidence decay:} How long do derived quantities remain valid?
  Should ADRs include a ``last reviewed'' date and a trigger for
  re-evaluation?
\item
  \textbf{Cross-ADR dependencies:} When multiple ADRs interact (e.g., a
  quantization choice interacts with a retrieval chunk-size choice), how
  should the team document the dependency graph?
\item
  \textbf{Tooling gaps:} No widely adopted tool exists for ADR graph
  visualization or impact analysis. Building such a tool would require
  understanding the ADR network and its evolution over time.
\item
  \textbf{ADR format evolution:} The minimal format (title, context,
  decision, consequences, alternatives) works for most cases, but edge
  cases --- decisions involving regulatory compliance, security, or data
  governance --- may require additional sections.
\end{enumerate}

These open questions are not blockers; they are signals for when the
practice matures and the team needs more sophisticated governance.

\section{End-of-Chapter Mini-Case}\label{end-of-chapter-mini-case}

\textbf{Scenario:} The fleet decides to migrate from a single 70B FP16
model to a two-model mixture-of-experts (MoE) architecture, each expert
34B parameters, FP16, running on the same A100 GPUs.

\textbf{The ADR (draft):}

\begin{itemize}
\tightlist
\item
  \textbf{Title:} ADR 0017 --- MoE Model Deployment
\item
  \textbf{Status:} Proposed
\item
  \textbf{Context:} User queries have grown 35\% YoY; a single 70B model
  cannot sustain p99 latency \textless{} 5 seconds under peak load. MoE
  allows routing 2× more effective parameters within the same memory
  budget.
\item
  \textbf{Decision:} Migrate to two 34B FP16 experts with top-1 routing,
  accepting a 5\% increase in inference latency per request due to
  routing overhead.
\item
  \textbf{Consequences:}

  \begin{itemize}
  \tightlist
  \item
    Positive: Effective parameter count rises from 70B to
    \textasciitilde68B (2 × 34B × routing fraction), improving answer
    quality on factual Q\&A. GPU memory per expert fits within A100 80
    GB, enabling 2× concurrent instances.
  \item
    Negative: Routing logic adds \textasciitilde50 ms latency; requires
    new monitoring for route distribution skew.
  \end{itemize}
\item
  \textbf{Alternatives considered:}

  \begin{enumerate}
  \def\labelenumi{\arabic{enumi}.}
  \tightlist
  \item
    Increase batch size --- rejected, increases memory pressure and
    worsens tail latency.
  \item
    Move to GGUF 4-bit --- rejected, quality regression above the SLA
    threshold.
  \item
    Model sharding across GPU nodes --- rejected, operations overhead
    for inter-node communication.
  \end{enumerate}
\item
  \textbf{Evidence:} FACT: A100 FP16 expert 34B footprint ≈ 68 GB
  parameters + 15 GB activations. DERIVED: Top-1 routing reduces
  effective parameters by \textasciitilde50\% compared to uniform
  mixing.
\end{itemize}

This mini-case illustrates how the ADR format captures not just the
``what'' but the ``how much'' --- derived quantities, trade-offs, and
explicit alternatives. The team can now evaluate the proposal against
the fleet's latency and quality SLAs, with all reasoning preserved for
future review.

\begin{center}\rule{0.5\linewidth}{0.5pt}\end{center}
